\appendix
\section{Formal Conditions for the Failure of Probabilistic Factorization}
  \label{app:factorization}

  This appendix provides a formal clarification of the structural conditions under
  which probabilistic factorization fails in the presence of non-injective projections.
  The purpose is not to introduce new physical assumptions, but to make explicit the
  logical content underlying the arguments presented in the main text.

  Let $(\Omega,\Sigma,\mu)$ be a probability space describing underlying configurations
  $\omega \in \Omega$, equipped with a normalized measure $\mu$.
  Observable outcomes are obtained through a measurable mapping
  \begin{equation}
    \Pi : \Omega \rightarrow \mathcal{O},
  \end{equation}
  where $\mathcal{O}$ denotes the space of observable descriptions.
  We assume that $\Pi$ is non-injective, so that there exist distinct
  $\omega_1 \neq \omega_2$ such that $\Pi(\omega_1)=\Pi(\omega_2)$.

  Let $A$ and $B$ denote two spacelike separated measurement settings, and let
  $a \in \mathcal{O}_A$ and $b \in \mathcal{O}_B$ denote the corresponding observable
  outcomes.
  The joint probability distribution accessible at the effective level is given by
  \begin{equation}
    P(a,b \mid A,B)
    =
    \mu\!\left(\Pi^{-1}(a,b)\right)
    =
    \sum_{\omega \in \Pi^{-1}(a,b)} \mu(\omega).
  \end{equation}

  Suppose that a factorizable representation of the joint distribution exists.
  Then there must exist a measurable variable $\lambda$ and a probability density
  $\rho(\lambda)$ such that
  \begin{equation}
    P(a,b \mid A,B)
    =
    \int d\lambda\,
    \rho(\lambda)\,
    P(a \mid A,\lambda)\,
    P(b \mid B,\lambda).
    \label{eq:bell-factorization}
  \end{equation}
  This representation implicitly assumes that $\lambda$ provides a complete and
  non-redundant parameterization of the relevant underlying configurations for the
  observable outcomes.

  However, when $\Pi$ is non-injective, no such variable $\lambda$ can generically
  exist at the effective level.
  Any candidate $\lambda$ constructed from observable data must identify entire
  equivalence classes
  \begin{equation}
    [\omega] = \{\omega' \in \Omega \mid \Pi(\omega')=\Pi(\omega)\},
  \end{equation}
  rather than individual underlying configurations.
  As a result, conditioning on $\lambda$ does not isolate independent contributions
  associated with $A$ and $B$.

  More formally, if $\lambda$ is a function of $\Pi(\omega)$, then $\lambda$ is constant
  over each equivalence class $[\omega]$.
  The joint probability conditioned on $\lambda$ therefore takes the form
  \begin{equation}
    P(a,b \mid A,B,\lambda)
    =
    \sum_{\omega \in [\omega_\lambda]} \mu(\omega \mid \lambda),
  \end{equation}
  which does not decompose into a product of independent marginal terms unless
  additional constraints are imposed.
  In generic non-injective cases, such constraints are absent, and factorization fails.

  This establishes that non-injectivity of the mapping from underlying configurations
  to observable descriptions is sufficient to obstruct any representation of the form
  Eq.~\eqref{eq:bell-factorization}.
  The failure of Bell-type factorization thus follows from a purely structural property
  of the descriptive mapping, independently of any assumptions about dynamical
  nonlocality, retrocausality, or hidden variables.

  The result emphasizes that Bell inequalities apply only to effective descriptions
  admitting an injective or equivalently factorizable representation.
  When observable outcomes arise through non-injective projections, violations of Bell
  inequalities are a natural and generic consequence of the descriptive structure itself.
